% 第2章 2.2.2：Riesz 表示的几何示意（实平面中的二维截面）。
% 核 N=ker L 为过原点的直线，表示向量 h_L 与 N 正交；水平集 {L=c} 是平行于 N 的直线，
% 相邻两条之间沿 h_L 方向相距 1/‖h_L‖。f 分解为 N 中分量与沿 h_L 的分量 (L(f)/‖h_L‖^2)h_L。
% 旋转坐标系中 h_L=(1.4,0)，f=(2.0,1.2)，比例 1 单位 = 10mm，整体旋转 25 度。
\begin{tikzpicture}[x=10mm,y=10mm,font=\figurefont\footnotesize,text=FigureInk,
  note/.style={font=\FigureNoteFont\fontsize{7.5}{9}\selectfont,text=FigureMuted},
  vec/.style={line width=1pt,-{Stealth[length=2mm]}}]
\begin{scope}[rotate=25]
  \foreach \c in {-1.4,1.4,2.8}{\draw[FigureGrid,line width=.8pt] (\c,-2.0)--(\c,2.3);}
  \draw[FigureBlue,line width=1.3pt] (0,-2.0)--(0,2.3);
  \draw[densely dotted,FigureYellow,line width=.8pt] (1.4,0)--(2.0,0);
  \draw[vec,FigureYellow] (0,0)--(1.4,0);
  \draw[vec,FigureInk] (0,0)--(2.0,1.2);
  \draw[densely dashed,FigureStroke,line width=.6pt] (2.0,1.2)--(0,1.2);
  \draw[densely dashed,FigureStroke,line width=.6pt] (2.0,1.2)--(2.0,0);
  \draw[vec,FigureGreen] (0,0)--(0,1.2);
  \draw[FigureStroke,line width=.5pt] (0.18,0)--(0.18,0.18)--(0,0.18);
  \coordinate (Ntop) at (0,2.3); \coordinate (L1) at (1.4,2.3); \coordinate (L2) at (2.8,2.3); \coordinate (Lm) at (-1.4,2.3);
  \coordinate (h) at (1.4,0); \coordinate (f) at (2.0,1.2); \coordinate (n) at (0,1.2); \coordinate (p) at (2.0,0);
\end{scope}
\fill[FigureInk] (0,0) circle (.6pt);
\node[text=FigureBlue,above] at (Ntop) {$N=\ker L$};
\node[note,above] at (L1) {$L=\|h_L\|^2$};
\node[note,above] at (L2) {$L=2\|h_L\|^2$};
\node[note,above] at (Lm) {$L=-\|h_L\|^2$};
\node[text=FigureYellow!70!FigureInk,below right=-.5mm] at (h) {$h_L$};
\node[right] at (f) {$f$};
\node[text=FigureGreen,left] at (n) {$f-\frac{L(f)}{\|h_L\|^2}h_L$};
\end{tikzpicture}
